\begin{table}[H]
\centering
\caption{Ranking departamental de PIB por trabajador y remuneraci\'on por trabajador, 2024pr}
\label{tab:pib_geih_productividad_departamento_remuneracion}
\scriptsize
\begin{tabular}{lrrrrr}
\toprule
Departamento & PIB/trab. 2024pr & Puesto & Rem./trab. 2024 & Puesto & Dif. puestos \\
\midrule
Bogotá D.C. & 67,4 & 1 & 2,95 & 1 & +0 \\
Meta & 63,1 & 2 & 1,83 & 4 & +2 \\
Santander & 62,0 & 3 & 1,80 & 6 & +3 \\
Valle del Cauca & 54,5 & 4 & 1,79 & 7 & +3 \\
Antioquia & 47,7 & 5 & 2,12 & 2 & -3 \\
Boyacá & 47,5 & 6 & 1,49 & 13 & +7 \\
Bolívar & 45,1 & 7 & 1,33 & 16 & +9 \\
Cundinamarca & 40,4 & 8 & 1,79 & 8 & +0 \\
Tolima & 40,0 & 9 & 1,50 & 11 & +2 \\
Risaralda & 38,7 & 10 & 1,82 & 5 & -5 \\
Atlántico & 37,3 & 11 & 1,58 & 10 & -1 \\
Caldas & 36,6 & 12 & 1,85 & 3 & -9 \\
Quindío & 33,8 & 13 & 1,71 & 9 & -4 \\
Cesar & 32,6 & 14 & 1,26 & 17 & +3 \\
Huila & 31,6 & 15 & 1,36 & 15 & +0 \\
Chocó & 28,6 & 16 & 1,16 & 18 & +2 \\
Cauca & 26,6 & 17 & 1,06 & 20 & +3 \\
Caquetá & 24,9 & 18 & 1,47 & 14 & -4 \\
La Guajira & 24,8 & 19 & 0,89 & 24 & +5 \\
Norte de Santander & 24,3 & 20 & 1,49 & 12 & -8 \\
Magdalena & 23,9 & 21 & 1,15 & 19 & -2 \\
Córdoba & 22,0 & 22 & 1,04 & 22 & +0 \\
Sucre & 22,0 & 23 & 1,01 & 23 & +0 \\
Nariño & 17,8 & 24 & 1,06 & 21 & -3 \\
\bottomrule
\end{tabular}
\caption*{\footnotesize Nota: PIB por trabajador en millones de pesos constantes de 2015; remuneraci\'on por trabajador en millones de pesos constantes de 2025 al mes. La remuneraci\'on por trabajador se calcula como la remuneraci\'on laboral mensual total observada entre ocupados con ingreso horario positivo y horas v\'alidas, dividida por el n\'umero total de ocupados del departamento. La diferencia de puestos corresponde al puesto en remuneraci\'on menos el puesto en PIB por trabajador; un valor positivo indica que el departamento ocupa una posici\'on m\'as baja en remuneraci\'on que en productividad. Fuente: c\'alculos propios con DANE y GEIH.}
\end{table}